
% Akademický úvod
\section{Úvod}

Tento text predstavuje systematický a didakticky usporiadaný prehľad problematiky Internetu vecí (Internet of Things, IoT) a vstavaných (embedded) systémov so zameraním na návrh, implementáciu a praktické overenie riešení. Skripta kombinuje akademicky podložené teoretické východiská s konkrétnymi praktickými príkladmi a experimentálnymi postupmi, ktoré sú vhodné pre použitie vo výučbe aj pri samostatnom štúdiu.

Cieľom je poskytnúť čitateľovi:
\begin{itemize}
	\item jasné pochopenie základných konceptov a architektúr vstavaných systémov;
	\item detailný prehľad komunikačných mechanizmov a protokolov bežne využívaných v IoT (vrátane ich vlastností, obmedzení a návrhových kompromisov);
	\item praktické návody pre implementáciu a ladenie riešení na reálnych platformách (s dôrazom na platformu ESP32 a prostredie MicroPython);
	\item súbor príkladov a cvičení umožňujúcich samostatné overenie a rozšírenie poznatkov.
\end{itemize}

Skriptum je koncipované pre študentov technických a prírodovedných programov, ako aj pre vývojárov prototypov a inžinierov, ktorí chcú získať systematický prehľad o metódach návrhu a praktických aspektoch nasadzovania IoT riešení. Predpokladá sa základná znalosť programovania (najmä jazykov Python alebo C) a elementárne znalosti elektroniky; tam, kde sú potrebné, sú základné pojmy kratko rekapitulované.

Metodologicky sa skripta opiera o kombináciu nasledovných prístupov: formálne definície a kategorizácie, komparatívna analýza alternatívnych technológií, návrhové vzory pri implementácii a praktické príklady s overiteľnými výsledkami. Autori kladú dôraz na to, aby boli vysvetlené dôvody návrhových rozhodnutí, nie iba opis postupov.

Organizácia textu: po úvodnom prehľade nasledujú kapitoly rozvíjajúce problematiku postupne do väčšej hĺbky. Najprv sú predstavené architektúry vstavaných systémov a základné hardvérové komponenty, potom sú rozpracované komunikačné topológie a protokoly, na ktoré nadväzujú praktické kapitoly venované konkrétnej platforme ESP32 a programovaniu v MicroPython. Záver tvorí súbor príkladov a cvičení s komentovanými riešeniami.

V rámci textu sú explicitne uvedené predpoklady a rozsah jednotlivých kapitol. Niektoré témy (napr. hlboká kybernetická bezpečnosť IoT, pokročilé cloudové integrácie alebo optimalizácia na úrovni architektúr SoC) sú zmienené ako odporučená literatúra a predmet ďalšieho štúdia, avšak nie sú v tomto skripte rozpracované do úplnej hĺbky.

Použitá terminológia a konvencie: tam, kde je to vhodné, používame formálne definície a jednotné označenie veličín. V kódoch a príkladoch sú uvedené jasné komentáre a referenciačné poznámky. Všetky zdroje a literatúra sú uvedené v bibliografii.

Záverom autori dúfajú, že táto skripta poslúži ako robustný východiskový materiál pre ďalšie akademické i praktické aktivity v oblasti IoT a vstavaných systémov.
